\subsection{NUTS有限定位：采样、诊断与内存保护}\label{sec:nuts-localization}
采用独立身份\texttt{windows-nuts-localization-v1}，沿用原生Windows CPU环境与原目标、初值、每链实际Python/NumPy/torch RNG起始状态。G1/L1/L2/H1/H2/M1各选择一个既有4096步失败输入，G2/A1/W1各选择一个既有有效输入；它们是预定机制案例，不是随机失败率样本。比较1/4个worker及旧\texttt{mcmc.diagnostics()}开启/整体旁路，每次始终四链、1024步预热、4096步采样、最大树深8、目标接受率0.8、对角质量适配及每worker一线程。worker数也改变进程复用，插桩持久化改变对象生命周期与I/O，因此不将新计时直接替换旧性能结果。

21项便携检查和3项原生Job检查的回执均通过，0失败、0跳过；独立二维高斯的四条件资格通过后才冻结主体。9输入共36次主体调用全部取得终态：34次完成、2次资源停止。按冻结模型/条件优先级，另选择G2/A1的四worker诊断开/关，共4次确认，2次完成、2次资源停止。含4次资格，共44次注册四链调用，低于48次上限；不重试同一调用身份、不替换失败输入。全部技术数组不计入正式后验样本或独立重复。

\begin{table}[htbp]\centering\small
\begin{tabular}{llcccc}\toprule
案例 & 原任务状态 & 1w开启 & 4w开启 & 1w旁路 & 4w旁路\\\midrule
@MATRIX@
\bottomrule\end{tabular}
\caption{九个固定案例的四条件主体结果。w表示worker数；完成仅指该技术工作流返回。G2/A1资源停止仍保留完整采样与失败状态。}
\end{table}

G2/A1主体及确认的诊断开启失败中，四链各1024步预热和4096步采样均已保存，采样钩子分块与诊断前完整数组逐位相同；末事件均为\texttt{legacy\_diagnostics\_enter}，没有退出标记，且收到原生Job内存限制消息。整体旁路均完成。Mac独立核对九模型54对条件的864项共同路径/RNG记录，以及64项确认/原条件记录，全部一致。六个选定历史失败案例在新插桩下四条件均完成；历史105项未分类失败、失败分母和成本保持不变。

独立诊断重放使用G1/G2/W1各一份保存的旧四链数组：旧单链ESS分别以1/4个worker组织，共6作业；R现代诊断对每份完整四链运行一次，共3作业。九作业中八次完成，G2四worker旧ESS在四链进入ESS后触发8 GiB集合RSS护栏。三组R联合诊断的二进制往返一致；Mac以R 4.6.0、posterior 1.7.0重算，Windows为R 4.6.1、posterior 1.7.0。74个参数的@R_NUMERIC_DIFFERENCES@个数值字段不同，最大按$\max(1,|a|,|b|)$缩放的差为$@R_MAX_SCALED_DIFFERENCE@$，未定位置及$\hat R>1.01$判断不变。

\begin{table}[htbp]\centering\small
\begin{tabular}{lrrr}\toprule
停止调用（均4w） & RSS观察峰/GiB & Job峰计数/GiB & 内存消息数\\\midrule
@RESOURCES@
\bottomrule\end{tabular}
\caption{全部五次资源停止的原始计数投影。RSS为进程工作集之和，可重复计入共享页且可能漏掉瞬时峰；Job原生峰计数与RSS、显存分别解释。护栏为8 GiB轮询RSS及12 GiB查询Job限制。}
\end{table}

RSS每0.2秒观察；私有内存只涵盖已捕获句柄且成功读取的部分，不当作全程精确峰。Job峰计数可超过查询限制，不能据此称硬保护失效或把它当作成功提交内存的精确峰；限制消息是直接证据，没有消息不能排除事件。部分失败退出码缺失，保持缺失。每次继续前核对原Job已终止和资源恢复，控制阈值及冻结执行源码未改。

冻结Pyro 1.9.2源码中，ESS的正配对序列长度为$4096/2-1=2047$，\texttt{\_cummin}将其repeat为$2047\times2047\times64$。float64重复数组为2145387008字节，扩展布尔掩码为268173376字节/链，另有FFT、输入和运行时内存。这里核对的是保存源码给出的静态规模；失败栈没有直接捕获这笔具体分配，不能把它写成实测独占分配峰。独立ESS重放与诊断进入后的内存消息共同支持该路径造成内存压力的解释，仍未建立全部历史故障的唯一原因。

Mac接收核对完整tar的30044个内容成员及149个执行文件的Git对象，确认选择由全部36项主体阶段结果独立重算。六份数值分析文件逐字节一致；README仅CRLF/LF不同，其校验清单相应不同，原件不改写。两次Windows终态只核验恢复均新增执行0，29444份当时研究文件集合/大小/哈希由Mac再次核对；保存的53个Job终态查询均为无活动。Mac核验保存证据，不声称查询了远端当前内核。允许的结论是本轮选定配置的诊断阶段资源故障和未复现案例，不是总体失败率、收敛或性能基准。
